% Part: first-order-logic
% Chapter: syntax-and-semantics
% Section: first-order-languages

\documentclass[../../../include/open-logic-section]{subfiles}

\begin{document}

\olfileid[fr]{fol}{syn}{fol}

\olsection{Langages du premier ordre}

Les expressions de la logique du premier ordre sont construites à
partir d'un vocabulaire de base comprenant des \emph{!!{variable}s},
des \emph{!!{constant}s}, des \emph{!!{predicate}s} et parfois des
\emph{!!{function}s}. Avec les connecteurs logiques, les
quantificateurs et les signes de ponctuation, comme les parenthèses
et les virgules, ces symboles permettent de former des \emph{termes}
et des \emph{!!{formula}s}.

\begin{explain}
De façon informelle, les !!{predicate}s désignent des propriétés et
des relations, les !!{constant}s des objets individuels, et les
!!{function}s des applications. À l'exception de l'!!{identity}~$\eq$,
ces symboles sont les \emph{symboles non logiques} ; leur ensemble
constitue un langage. Tout langage du premier ordre~$\Lang L$ est
déterminé par ses symboles non logiques. Dans le cas le plus général,
$\Lang L$ contient une infinité de symboles de chaque sorte.
\end{explain}

Dans le cas général, nous utilisons les symboles suivants en logique
du premier ordre :

\begin{enumerate}
\item Symboles logiques
\begin{enumerate}
\item Connecteurs logiques et quantificateurs :
  \startycommalist
  \iftag{prvNot}{\ycomma $\lnot$ (négation)}{}%
  \iftag{prvAnd}{\ycomma $\land$ (conjonction)}{}%
  \iftag{prvOr}{\ycomma $\lor$ (disjonction)}{}%
  \iftag{prvIf}{\ycomma $\lif$ (!!{conditional})}{}%
  \iftag{prvIff}{\ycomma $\liff$ (!!{biconditional})}{}%
  \iftag{prvAll}{\ycomma $\lforall$ (quantificateur universel)}{}%
  \iftag{prvEx}{\ycomma $\lexists$ (quantificateur existentiel)}{}.
\tagitem{prvFalse}{La constante propositionnelle de
  !!{falsity}~$\lfalse$.}{}
\tagitem{prvTrue}{La constante propositionnelle de
  !!{truth}~$\ltrue$.}{}
\item L'!!{identity} binaire~$\eq$.
\item Un ensemble !!{denumerable} de !!{variable}s : $\Obj v_0$,
  $\Obj v_1$, $\Obj v_2$, \dots
\end{enumerate}
\item Symboles non logiques, qui constituent le \emph{langage standard}
  de la logique du premier ordre
\begin{enumerate}
\item Un ensemble !!{denumerable} de !!{predicate}s à $n$ places
  pour chaque $n>0$ : $\Obj A^n_0$, $\Obj A^n_1$,
  $\Obj A^n_2$, \dots
\item Un ensemble !!{denumerable} de !!{constant}s : $\Obj c_0$,
  $\Obj c_1$, $\Obj c_2$, \dots.
\item Un ensemble !!{denumerable} de !!{function}s à $n$ places
  pour chaque $n>0$ : $\Obj f^n_0$, $\Obj f^n_1$,
  $\Obj f^n_2$, \dots
\end{enumerate}
\item Signes de ponctuation : (, ) et la virgule.
\end{enumerate}

La plupart de nos définitions et de nos résultats porteront sur le
langage standard complet de la logique du premier ordre. Selon
l'application envisagée, nous pourrons toutefois restreindre le
langage à quelques !!{predicate}s, !!{constant}s et !!{function}s.

\begin{ex}
Le langage~$\Lang L_A$ de l'arithmétique contient un unique
!!{predicate} à deux places~$<$, une unique !!{constant}~$\Obj 0$,
un !!{function} à une place~$\prime$, et deux !!{function}s à deux
places, $+$ et~$\times$.
\end{ex}

\begin{ex}
Le langage~$\Lang L_Z$ de la théorie des ensembles ne contient qu'un
seul !!{predicate} à deux places~$\in$.
\end{ex}

\begin{ex}
Le langage~$\Lang L_\le$ des ordres ne contient que le
!!{predicate} à deux places~$\le$.
\end{ex}

Il s'agit, là encore, de conventions : ces symboles ne sont
officiellement que des alias. Ainsi, $<$, $\in$ et $\le$ sont des
alias de $\Obj A^2_0$, $\Obj 0$ de $\Obj c_0$, $\prime$ de
$\Obj f^1_0$, $+$ de $\Obj f^2_0$ et $\times$ de $\Obj f^2_1$.

\iftag{defNot,defOr,defAnd,defIf,defIff,defTrue,defFalse,defEx,defAll}{%
Outre les connecteurs primitifs et
\iftag{notprvEx,notprvAll}{le quantificateur}{les quantificateurs}
introduits ci-dessus, nous utilisons aussi les symboles
\emph{définis} suivants :
\startycommalist
  \iftag{defNot}{\ycomma $\lnot$ (négation)}{}%
  \iftag{defAnd}{\ycomma $\land$ (conjonction)}{}%
  \iftag{defOr}{\ycomma $\lor$ (disjonction)}{}%
  \iftag{defIf}{\ycomma $\lif$ (!!{conditional})}{}%
  \iftag{defIff}{\ycomma $\liff$ (!!{biconditional})}{}%
  \iftag{defAll}{\ycomma $\lforall$ (quantificateur universel)}{}%
  \iftag{defEx}{\ycomma $\lexists$ (quantificateur existentiel)}{}%
  \iftag{defFalse}{\ycomma !!{falsity}~$\lfalse$}{}%
  \iftag{defTrue}{\ycomma !!{truth}~$\ltrue$}}{}.

\begin{tagblock}{defNot,defOr,defAnd,defIf,defIff,defTrue,defFalse,defEx,defAll}
\begin{explain}
Un symbole défini ne fait pas officiellement partie du langage : il
est introduit comme une abréviation informelle. Il permet de raccourcir
des !!{formula}s qui, écrites avec les seuls symboles primitifs,
seraient assez longues. L'avantage est évident ; plus encore, les
démonstrations deviennent plus courtes. Tout symbole primitif doit en
effet être traité séparément dans les démonstrations. Plus il y a de
symboles primitifs, plus celles-ci sont longues.
\end{explain}
\end{tagblock}

% Autres symboles

\begin{intro}
Vous connaissez peut-être une terminologie ou des symboles différents
de ceux employés ci-dessus. Les ouvrages et les enseignants en logique
utilisent couramment $\sim$, $\neg$ ou~!\ pour la « négation », et
$\wedge$, $\cdot$ ou $\&$ pour la « conjonction ». Pour le
« conditionnel » ou l'« implication », on rencontre souvent
$\rightarrow$, $\Rightarrow$ et $\supset$.
\iftag{prvIff,defIff}{Les symboles du « biconditionnel », de la
  « bi-implication » ou de l'« équivalence (matérielle) » sont
  $\leftrightarrow$, $\Leftrightarrow$ et $\equiv$.}{}
\iftag{prvFalse,defFalse}{Le symbole $\lfalse$ est appelé, selon les
  auteurs, « fausseté », « falsum », « absurdité » ou « bottom ».}{}
\iftag{prvTrue,defTrue}{Le symbole $\ltrue$ est appelé « vérité »,
  « verum » ou « top ».}{}

On emploie habituellement des lettres minuscules du début de
l'alphabet latin, par exemple $a$, $b$, $c$, pour les !!{constant}s,
parfois appelées noms, et des lettres minuscules de la fin de
l'alphabet, par exemple $x$, $y$, $z$, pour les !!{variable}s. Les
quantificateurs se combinent avec des !!{variable}s, comme~$x$ ;
pour le quantificateur universel, on trouve notamment les notations
$\forall x$, $(\forall x)$, $(x)$, $\Pi x$ et $\bigwedge_x$ ; pour
le quantificateur existentiel, $\exists x$, $(\exists x)$, $(Ex)$,
$\Sigma x$ et $\bigvee_x$.
\end{intro}

\begin{explain}
Nous pourrions prendre tous les !!{operator}s propositionnels et les
deux quantificateurs comme symboles primitifs du langage. Nous
pourrions aussi choisir un ensemble plus réduit de symboles primitifs
et traiter les autres !!{operator}s comme des symboles définis. Parmi
les ensembles de connecteurs booléens fonctionnellement complets
figurent $\{ \lnot, \lor \}$, $\{ \lnot, \land \}$ et
$\{ \lnot, \lif\}$ ; chacun peut être combiné avec l'un ou l'autre
des quantificateurs pour obtenir un langage du premier ordre d'une
expressivité complète.

Vous connaissez peut-être deux autres !!{operator}s : la barre de
Sheffer~$|$, nommée d'après Henry Sheffer, et la flèche de
Peirce~$\downarrow$, également appelée « dagger de Quine ». Dans
leurs interprétations usuelles respectives, « non-et » (nand) et
« non-ou » (nor), chacun de ces opérateurs est à lui seul
fonctionnellement complet.
\end{explain}

\end{document}
